% =========================================================================
%  Capítulo 1 — Introdução
% =========================================================================

\chapter{Introdução}
\label{cap:introducao}

% --- INSTRUÇÕES (apagar antes da entrega) ---------------------
% A introdução deve permitir a um leitor externo perceber:
%   - O problema que o sistema vem resolver.
%   - O contexto em que o sistema vai ser utilizado.
%   - Quais os objetivos do projeto.
%   - Como está organizado o restante do relatório.
% ---------------------------------------------------------------

\section{Enquadramento}
\label{sec:enquadramento}

% Apresentem aqui o domínio do problema. O que é a GREENHERB?
% Que tipo de operação se passa numa estufa de ervas aromáticas?
% Que dificuldades existem hoje sem informatização?
% (1 a 2 parágrafos)

A empresa GREENHERB pretende informatizar a gestão de uma estufa
dedicada à produção de ervas aromáticas, designadamente
manjericão, hortelã, alecrim e tomilho. O sistema atual,
maioritariamente manual, cria dificuldades no planeamento, na
monitorização das condições ambientais e na rastreabilidade das
intervenções realizadas.

[Completar: descrever o impacto destas dificuldades na operação
quotidiana da estufa e a oportunidade que a digitalização representa.]


\section{Objetivos}
\label{sec:objetivos}

% Listar os objetivos do projeto. Recomenda-se uma lista
% numerada, com objetivos verificáveis (não vagos).

O presente trabalho tem como objetivos:

\begin{enumerate}
  \item Desenvolver uma aplicação web \emph{full-stack} que suporte o ciclo de vida completo de lotes de cultivo;
  \item Implementar autenticação, autorização baseada em perfis e rastreabilidade das ações realizadas;
  \item Especificar e documentar a API segundo a norma OpenAPI 3.x;
  \item Definir uma arquitetura de armazenamento no browser que permita o funcionamento da aplicação em condições de conectividade instável;
  \item {}[Adicionar / ajustar conforme decisões da equipa.]
\end{enumerate}


\section{Contexto e restrições}
\label{sec:contexto}

% Esta secção é onde se reconhecem as restrições do projeto.
% Útil para enquadrar decisões posteriores.

A solução é desenvolvida com restrições técnicas previamente
definidas: o \emph{frontend} em \texttt{HTML}, \texttt{CSS} e
\texttt{JavaScript}, o \emph{backend} em \texttt{Node.js} com
recurso a \texttt{Mongoose} para acesso a uma base de dados
\texttt{MongoDB}, autenticação baseada em \texttt{JWT} e uso obrigatório
dos mecanismos de armazenamento do navegador (\texttt{localStorage}
ou \texttt{sessionStorage}, \texttt{IndexedDB} e \texttt{Cache API}).

Acresce o contexto operacional: a aplicação é utilizada em ambiente
de estufa onde a conectividade pode ser instável, o que condiciona
fortemente as decisões de arquitetura do lado do cliente.


\section{Estrutura do documento}
\label{sec:estrutura}

% Pequena ``mapa'' do que o leitor vai encontrar nos próximos
% capítulos. Manter curto.

O restante documento organiza-se da seguinte forma. O Capítulo~\ref{cap:requisitos}
identifica e especifica os requisitos do sistema, distinguindo entre
funcionais e não funcionais e apresentando a sua priorização. O
Capítulo~\ref{cap:arquitetura} descreve a arquitetura adotada,
detalhando as responsabilidades de cada camada (\emph{browser}, API e
base de dados) e justificando as decisões tomadas. O
Capítulo~\ref{cap:conclusao} apresenta as conclusões, balanço crítico
e trabalho futuro.
